\begin{scanpage}{0015}{xi}
wave which rapidly decreases with the distance from the edge. Therefore, one may consider that the nonuniform part of the current is concentrated mainly in the vicinity of the discontinuity. By means of this physically obvious assumption, the field scattered by a strip (Chapter I), by a disk (Chapter II), by a finite length cylinder (Chapter III) and by certain other bodies of rotation (Chapter IV) is calculated.

For a more precise calculation, however, it is necessary to keep in mind that the actual current distribution in the vicinity of the body's edges differs from the current distribution near the edge of the wedge. Actually, the edge wave corresponding to the nonuniform part of the current, propagated along the body's surface, reaches the adjacent edge and undergoes diffraction by it, exciting secondary edge waves. The latter in turn produce new edge waves, etc. If all the dimensions of the body are large in comparison with the wavelength, then as a rule it is sufficient to consider only the secondary diffraction. This phenomenon is studied in Chapter V using the example of a strip and disk.

In the case of a narrow cylindrical conductor of finite length, the edge waves of the current decrease very slowly with the distance from each end. Therefore, here it is impossible to limit oneself to a consideration only of secondary diffraction, and it is necessary to investigate the multiple diffraction of edge waves. Chapter VII is devoted to this problem.

The uniform and nonuniform parts of the current are more than auxiliary concepts which are useful in solving diffraction problems. In Chapter VI it is shown that one is able experimentally to separate from the total fringing field that part of it which is created by the nonuniform part of the current. There, it is also shown that the depolarization phenomenon of the reflected signal is caused \uline{only} by the nonuniform part of the current.
\end{scanpage}
